\subsection{数量微分方程线性组的基本积分组}

在本小节及下一小节中，我们将考虑如下情形：E 是复数域 C 上的 n 维向量空间（因而在 R 上是 2n 维的），并且对每个 \(t \in J\) ，\(A(t)\) 都是 E 关于 C 上向量空间结构的一个自同态。于是可以把 \(A(t)\) 与它关于 E 的一个基（在域 C 上）的矩阵 \((a_{ij}(t))\) 等同起来，这时的 \(a_{ij}\) 是 J 上有定义且规则的 \(n^{2}\) 个复值函数；若以 \(x_{j}\) ( \(1 \leqslant j \leqslant n\) ) 表示向量 \(x \in E\) 关于所选基的（复）分量，则线性方程

\[
{\frac {d \mathbf {x}}{d t}} = A (t). \mathbf {x} + \mathbf {b} (t)\tag{2}
\]

仍等价于方程组

\[
\frac {d x _ {i}}{d t} = \sum_ {j = 1} ^ {n} a _ {i j} (t) x _ {j} + b _ {i} (t) \quad (1 \leqslant i \leqslant n).\tag{3}
\]

于是，定理 1 (IV, p.~179) 与定理 2 (IV, p.~179) 以及命题 2 (IV, p.~181) 表明，对 E 内的每个 \(\mathbf{x}_0 = (x_{k0})_{1\leqslant k\leqslant n}\) ，存在一个且仅一个积分 \(\mathbf{u} = (u_k)_{1\leqslant k\leqslant n}\) 满足方程

\[
{\frac {d \mathbf {x}}{d t}} = A (t). \mathbf {x}\tag{4}
\]

它在 E 上定义，并且取值 \(x_{0}\) 于点 \(t_{0}\) ；这个积分可以写成

\[
\mathbf {u} (t, t _ {0}, \mathbf {x} _ {0}) = C (t, t _ {0}). \mathbf {x} _ {0},
\]

其中 \(C(t, t_{0})\) 是一个 n 阶可逆方阵 \((c_{ij}(t, t_{0}))\) ，其元素是 J × J 上的连续复值函数，并且使得 \(t \mapsto c_{ij}(t, t_{0})\) 是 J 上某个规则函数的原函数。

在 \(n = 1\) 的特殊情形，方程组 (3) 化为单个数量方程

\[
{\frac {d x}{d t}} = a (t) x + b (t)\tag{11}
\]

\((a(t) \text{ 与 } b(t) \text{ 为 J 上的复值规则函数}); \text{ 立即可验证，(单元素)矩阵 } C(t, t_0) \text{ 等于 } \exp\left(\int_{t_0}^{t} a(s) ds\right); \text{ (11) 的积分，其值等于 } x_0 \text{，在点 } t_0 \text{ 处，因而明确地由下式给出}\)

\[
u (t) = x _ {0} \exp \left(\int_ {t _ {0}} ^ {t} a (s) d s\right) + \int_ {t _ {0}} ^ {t} b (s) \exp \left(\int_ {t _ {0}} ^ {t} a (\tau) d \tau\right) d s.\tag{12}
\]

在由 J 到 E 的连续映射组成、赋予紧收敛拓扑的空间 \(\mathcal{C}(\mathrm{J};\mathrm{E})\) 中，方程 (4) 的积分所成的集 I 是一个向量子空间（在 C 上），它同构于 E，从而同构于 \(C^{n}\) (IV, p.~180, cor. 1 及 IV, p.~181, prop. 2)。

这个空间（在域 C 上）的一个基 \((\mathbf{u}_{j})_{1\leqslant j\leqslant n}\) 称为 (4) 的一个基本积分组。

命题 4. 为使方程 (4) 的 n 个积分 \(\mathbf{u}_{j}\ (1 \leqslant j \leqslant n)\) 构成一个基本积分组，必须且只需它们的值 \(\mathbf{u}_{j}(t_{0})\) （在某个点 \(t_{0} \in J\) 处取值）是 E 中线性无关的向量。

事实上，把每个 \(\mathbf{x}_0 \in \mathrm{E}\) 映为积分 \(t \mapsto C(t, t_0). \mathbf{x}_0\) 的映射是 \(\mathrm{E}\) 到 \(\mathcal{I}\) 上的一个同构 (IV, p.~180, cor. 1 及 IV, p.~181, prop. 2)。

若 \(\left(\mathbf{e}_j\right)_{1\leqslant j\leqslant n}\) 是 E 在 C 上的任意一个基，则 \(n\) 个积分

\[
\mathbf {u} _ {j} (t) = C \left(t, t _ {0}\right). \mathbf {e} _ {j} \quad (1 \leqslant j \leqslant n)
\]

便构成一个基本积分组；若把 \(C(t, t_0)\) 与它关于基 \((\mathbf{e}_j)\) 的矩阵等同起来，则诸积分 \(\mathbf{u}_j\) 正是矩阵 \(C(t, t_0)\) 的各列。于是，(4) 的取值 \(\mathbf{x}_0 = \sum_{j=1}^{n} \lambda_j \mathbf{e}_j\) 于点 \(t_0\) 的积分就是 \(C(t, t_0). \mathbf{x}_0 = \sum_{k=1}^{n} \lambda_k \mathbf{u}_k(t)\) 。

对 (4) 的任意 \(n\) 个积分 \(\mathbf{u}_j\) ( \(1 \leqslant j \leqslant n\) )，我们把这 \(n\) 个积分在点 \(t \in \mathbf{J}\) 处关于基 \((\mathbf{e}_j)_{1 \leqslant j \leqslant n}\) （在 \(E\) 中）的行列式称为

\[
\Delta (t) = \left(\mathbf {u} _ {1} (t), \mathbf {u} _ {2} (t), \dots , \mathbf {u} _ {n} (t)\right)\tag{13}
\]

即 \(n\) 个向量 \(\mathbf{u}_j(t)\) 关于基 \((\mathbf{e}_j)\) 的行列式 (Alg., III, p.~522)。我们有 (Alg., III, p.~523, prop. 2)

\[
\Delta (t) = \Delta (t _ {0}) \det \left(C (t, t _ {0})\right).\tag{14}
\]

由 IV, p.~184 的命题 4，为使 \(\left(\mathbf{u}_{j}\right)_{1\leqslant j\leqslant n}\) 是 (4) 的一个基本积分组，必须且只需行列式 \(\Delta(t)\) （诸 \(u_{j}\) 的行列式）\(\neq0\) 在 J 的某个点 \(t_{0}\) 处成立；于是公式 (14) 表明，在 J 的每一点处 \(\Delta(t)\neq0\)，换言之，向量 \(\mathbf{u}_{j}(t)\) ( \(1\leqslant j\leqslant n\) ) 恒线性无关。

命题 5. 矩阵 \(C(t, t_{0})\) 的行列式由下列公式给出

\[
\det \left(C (t, t _ {0})\right) = \exp \left(\int_ {t _ {0}} ^ {t} \operatorname{Tr} (A (s)) d s\right).\tag{15}
\]

事实上，若令 \(\delta(t) = \det\left(C(t, t_0)\right)\) ，则由行列式的求导公式 (I, p.~8, 公式 (3)) 有

\[
\frac {d \delta}{d t} = \operatorname{Tr} \left(\frac {d C (t , t _ {0})}{d t} \left(C (t, t _ {0})\right) ^ {- 1}\right) \delta (t)
\]

亦即，由 \(C(t, t_0)\) 所满足的 IV, p.~179 的微分方程 (5)，有

\[
\frac {d \delta}{d t} = \mathrm{Tr} \big (A (t) \big) \delta (t).
\]

由于 \(\delta(t_0) = 1\) ，由数量线性方程的积分的表达式 (12) (IV, p.~183) 即得公式 (15)。

正如刚才所见，给定 (4) 的 n 个线性无关的积分，便确定了这个方程的全部积分。现在我们要证明，当 \(1 \leqslant p \leqslant n\) 时，知道了方程 (4) 的 p 个线性无关的积分 \(\mathbf{u}_{j}\) ( \(1 \leqslant j \leqslant p\) )，便可把这个方程的积分化为由 n - p 个数量方程组成的齐次线性方程组的积分。设在一个区间 \(K \subset J\) 上存在 n - p 个映射

\[
\mathbf {u} _ {p + k} \quad (1 \leqslant k \leqslant n - p)
\]

它们把 K 映入 E，都是 K 上规则函数的原函数，并且使得对每个 \(t \in K\) ，n 个向量 \(\mathbf{u}_{j}(t) (1 \leqslant j \leqslant n)\) 构成 E 的一个基。

对每个点 \(t_1 \in \mathbf{J}\) ，总存在一个区间 \(\mathbf{K}\) ，即 \(t_1\) 在 \(\mathbf{J}\) 中的一个邻域，在其上可以定义具有前述性质的 \(n - p\) 个函数 \(\mathbf{u}_{p+k}\) ( \(1 \leqslant k \leqslant n - p\) )。因为，设 \((\mathbf{e}_i)_{1 \leqslant i \leqslant n}\) 是 \(\mathbf{E}\) 的一个基；该基中存在 \(n - p\) 个向量，它们与诸 \(\mathbf{u}_j(t_1)\) ( \(1 \leqslant j \leqslant p\) ) 一起构成 \(\mathbf{E}\) 的一个基 (Alg., II, p.~292, th. 2)；例如设它们是 \(\mathbf{e}_{p+1}, \ldots, \mathbf{e}_n\) ；由于行列式 \(\det(\mathbf{u}_1(t), \ldots, \mathbf{u}_p(t), \mathbf{e}_{p+1}, \ldots, \mathbf{e}_n)\) （关于基 \((\mathbf{e}_i)\) ）是 \(t\) 的连续函数并且在 \(t = t_1\) 处不为零，故存在一个邻域 \(\mathbf{K}\) 包含 \(t_1\) ，使行列式在其上不为零；于是可取 \(\mathbf{u}_{p+k}(t) = \mathbf{e}_{p+k}\) ( \(1 \leqslant k \leqslant n - p\) ) （对 \(t \in \mathbf{K}\) ）。

存在一个 n 阶可逆矩阵 \(B(t)\) ，其元素是 K 上规则函数的原函数，使得 \(B(t).e_{j} = \mathbf{u}_{j}(t)\) （对 \(1 \leqslant j \leqslant n\) ）。令 \(\mathbf{x} = B(t).y\) ；则 y 满足方程 \(\frac{dB}{dt}.y + B(t).\frac{dy}{dt} = A(t)B(t).y\) ，它也可以写成

\[
\frac {d \mathbf {y}}{d t} = (B (t)) ^ {- 1} \left(A (t) B (t) - \frac {d B}{d t}\right). \mathbf {y} = H (t). \mathbf {y}
\]

其中 \(H(t) = \left(h_{jk}(t)\right)\) 是一个元素在 K 上为规则函数的矩阵。由 \(B(t)\) 的定义，这个线性方程以 p 个常向量 \(\mathbf{e}_{j}\) ( \(1 \leqslant j \leqslant p\) ) 为积分；由此立即推出，必有 \(h_{jk}(t) = 0\) （对 \(1 \leqslant k \leqslant p\) ）；于是 y 的分量 \(y_{k}\) （关于基 \((\mathbf{e}_{i})\) ）中下标 \(k \geqslant p + 1\) 的那些，满足一个由 n - p 个方程组成的齐次线性方程组；一旦确定了这个方程组的解，\(dy_{j}/dt\) （对下标 \(j \leqslant p\) ）便是 \(y_{k}\) （其中 \(k \geqslant p + 1\) ）的线性函数，从而是已知的，而这些函数的原函数便给出 \(y_{j}\) （对下标 \(j \leqslant p\) ）。

特别地，当已知 IV, p.~183 的方程 (4) 的 \(n - 1\) 个线性无关的积分时，这个方程的积分便化为单个齐次数量方程的积分，随后只需计算 \(n\) 个原函数。

注. 1) 上述一切也适用于如下情形：E 是域 R 上的 n 维空间，并且 \(A(t)\) 对每个 \(t \in J\) 都是 E 的一个自同态：只需处处把 C 换成 R。

\begin{enumerate}
\def\labelenumi{\arabic{enumi})}
\setcounter{enumi}{1}
\tightlist
\item
  设 \(A(t) = (a_{ij}(t))\) 是一个 n 阶方阵，其元素是 t 在 J 上的实值（相应地，复值）规则函数，又设 \(C(t, t_0) = (c_{ij}(t, t_0))\) 是相应的线性方程组 (3) (IV, p.~22) 的预解矩阵。设 F 是 R（相应地，C）上的任意一个完备赋范空间，并考虑线性微分方程组
\end{enumerate}

\[
\frac {d \mathbf {y} _ {i}}{d t} = \sum_ {j = 1} ^ {n} a _ {i j} (t) \mathbf {y} _ {j}
\]

其中未知函数 \(\mathbf{y}_j\) 取值于 F。立即可见，该方程组的解 \((\mathbf{u}_j)_{1\leqslant j\leqslant n}\) ——它满足 \(\mathbf{u}_j(t_0) = \mathbf{d}_j\) （对 \(1\leqslant j\leqslant n\) ； \(\mathbf{d}_j\) 在 F 中任意取值）——由下列公式给出

\[
\mathbf {u} _ {i} (t) = \sum_ {j = 1} ^ {n} c _ {i j} (t, t _ {0}) \mathbf {d} _ {j} \quad (1 \leqslant i \leqslant n).
\]

我们特别考虑如下情形：\(A(t)\) 是 C 上 n 维向量空间 E 的一个自同态，并且存在 E 的一个基，使 \(A(t)\) 关于该基的矩阵对每个 \(t \in J\) 都有实元素。于是上面的结果表明（由 IV, p.~179 的定理 1），预解矩阵 \(C(t, t_{0})\) 关于同一个基也有实元素：只需考虑由所考虑的 E 的基在 R 上生成的向量空间 \(E_{0}\) ，并注意到 \(A(t)\) 在 \(E_{0}\) 上的限制是这个向量空间的一个自同态。

